\documentclass[11pt,a4paper]{article}
\usepackage[utf8]{inputenc}
\usepackage[T1]{fontenc}
\usepackage[french]{babel}
\usepackage{amsmath,amssymb}
\usepackage{geometry}
\geometry{hmargin=2.5cm,vmargin=2.5cm}

\title{TP Informatique : Approximation de $\pi$}
\author{Correction}
\date{}

\begin{document}
\maketitle

\section{Approximation de $\pi$}

On considère les suites couplées $(u(k))_{k \in \mathbb{N}}$ et $(v(k))_{k \in \mathbb{N}}$ définies par :

$$
\begin{cases}
u(0) = 1 \\
v(0) = 2 \\
u(k+1) = \dfrac{2\,u(k)\,v(k)}{u(k) + v(k)} \\
v(k+1) = \sqrt{u(k+1)\,v(k)}
\end{cases}
$$

On admet que $v(k) \xrightarrow[k \to +\infty]{} \dfrac{\sqrt{27}}{\pi}$.

\subsection*{Question 1}
Écrire un programme renvoyant pour un entier $N$ l'approximation du nombre $\pi$ obtenue à partir de la valeur de $v(N)$.

\subsubsection*{Correction 1}
Pour calculer $v(N)$, on utilise une boucle bornée en mettant à jour simultanément ou séquentiellement les variables $u$ et $v$. Attention à bien calculer le nouveau $u$ avant d'évaluer le nouveau $v$.

\begin{verbatim}
from math import sqrt

def approx_pi_rang(N):
    u = 1
    v = 2
    for k in range(N):
        u = (2 * u * v) / (u + v)
        v = sqrt(u * v)
    return sqrt(27) / v
\end{verbatim}

\subsection*{Question 2}
Écrire un programme renvoyant pour un réel $\epsilon > 0$ l'approximation de $\pi$ obtenue à partir de la valeur $v(n_0)$, où $n_0$ est le premier entier tel que :

$$
\left| \frac{u(n_0) - v(n_0)}{u(n_0) + v(n_0)} \right| < \epsilon
$$

\subsubsection*{Correction 2}
On utilise une boucle non bornée (\texttt{while}) qui s'arrête dès que la condition d'écart relatif est satisfaite.

\begin{verbatim}
from math import sqrt

def approx_pi_precision(epsilon):
    u = 1
    v = 2
    while abs((u - v) / (u + v)) >= epsilon:
        u = (2 * u * v) / (u + v)
        v = sqrt(u * v)
    return sqrt(27) / v
\end{verbatim}

\end{document}